\section{Separating Extraction Reliability from Detection}
\label{sec:protocol}

We distinguish three observable questions: can the response be processed under
the requested contract, can its proposed items be located in the source, and
do the located spans match the benchmark policy? These questions share an
evaluation collection but have different denominators. A response may satisfy
the JSON contract while returning no PII, or contain several correct quotations
alongside one invalid item. Neither structural compliance nor strict extraction
F1 alone explains both cases.

\subsection{A fixed output task under different context}
Let a document contain $n$ Unicode code points. For request $j$, define a
non-overlapping core $C_j=[s_j,e_j)$ of at most $c=1{,}200$ characters and a
target $T_j=[s_j,\min(n,e_j+h))$, with following halo $h=1{,}200$.
The full condition supplies the entire document as context. The local
condition supplies $[\max(0,s_j-h),\min(n,e_j+h))$. Both also supply the
same target string and core length, with the same category definitions and
output cap. Partitioning uses text length only; it does not inspect gold
categories, subject identities or span locations.

An output item specifies a category, an exact target quotation and a
one-based occurrence index, counting overlapping occurrences. Its span is
owned by the request containing its start in $C_j$; its end may extend into
the following halo but not beyond $T_j$. This ownership rule avoids counting
the same mention from overlapping requests. Identical category--boundary
predictions are deduplicated deterministically. The scheme avoids asking the
model to compute document-wide numeric offsets, but exact quotation and
occurrence selection remain part of the task. A wrong occurrence can yield
an incorrect location even when the text exists elsewhere in the target.

The output task is held fixed, not the input length or inference cost.
Mentions extending beyond their owner's target remain gold false negatives
if they cannot be represented. We do not shorten the gold collection to make
either context condition easier. The experiment therefore measures the
effect of additional input context within this particular extraction
interface, rather than comparing unconstrained whole-document summarization
with chunk-by-chunk annotation.

\subsection{Primary strict extraction}
Strict evaluation requires normal completion, the specified JSON root and
item structure, known labels, exact quotation anchoring and valid ownership
for every item in a response. If any item fails, that request contributes
no predictions. Truncated, terminal-failure and unparseable responses are
also empty. Gold spans remain in every case. Micro-precision and recall use
corpus-aggregated exact typed-span counts, with
\begin{equation}
 F_1=\frac{2\mathrm{TP}}{2\mathrm{TP}+\mathrm{FP}+\mathrm{FN}}.
\end{equation}
This primary score measures end-to-end extraction under the frozen contract.
Atomic rejection can suppress both valid and invalid items in the same
response, so it should not be read as a direct count of all recognizable PII.

\subsection{Post-hoc diagnostic views}
We replay archived responses without another model call. The diagnostic
parser permits removal of a single whole-response JSON fence and validates
parsed items independently. It retains exactly anchored valid items and
charges one exact false positive for each invalid parsed item. Let $I$ be
this invalid-item count and $\mathrm{FP}_a$ the unmatched anchored predictions.
The detection proxy uses $\mathrm{FP}_D=\mathrm{FP}_a+I$ in precision and
F1; recall includes all gold, including gold belonging to failed requests.
Thus keeping valid siblings does not make malformed items disappear without
a penalty. We call these measures D-P, D-R and D-F1.

Structural compliance is measured separately over \emph{all requests}.
It requires normally completed raw JSON with the prescribed field types
and no duplicate keys; JSON fences fail this test. It does not require
that a category belongs to the taxonomy or that a quotation exists in the
target. The empty-list rate also uses all requests and is computed after
the diagnostic whole-response fence parsing. A well-formed empty list can
therefore increase compliance without recovering any sensitive text.

Grounding diagnostics classify parsed-item failures by schema, category,
quotation, occurrence and core ownership. Gold-oriented summaries further
distinguish exact hits, wrong labels at exact boundaries, partial overlaps
and unresolved misses. These priority categories describe observations;
they do not identify the cause of a model's error. An unanchorable item
cannot be treated as evidence that any particular gold mention was detected.

No view repairs spelling, boundaries or labels, searches fuzzily, consults
gold to reinterpret an output, or drops failed requests. Terminal and
unparseable responses contain no usable parsed-item count and remain empty;
we do not invent an item penalty for them. The diagnostic also preserves
the original JSON parser's last-value behavior for duplicate keys while
flagging them as noncompliant. It cannot recover overwritten values.
Consequently D-F1 is an observable detection proxy, not format-independent
latent ability, and neither D-F1 nor character overlap replaces the primary
strict score.
